%! @@TITLE@@ — Unit @@UNITINT@@, Lesson @@LESSONID@@ — Lesson Plan
% SLIDES-FIRST plan (course shape 2026-09-25): the deck carries the lesson — warm-up, then 3–4
% cycles of definition → worked example → Now you try, a wrap-up, and the homework started in
% class. Students take their notes on the printed slide handout. Teacher-facing; never handed out.
\documentclass[10pt]{article}
\usepackage{@@PREFIX@@-article}
\usepackage{@@PREFIX@@-boxes}
\usepackage{graphicx}   % -article does NOT load graphicx; needed for thumbnails/screenshots
\graphicspath{{images/}}
\newcommand{\TallMath}[1]{$\displaystyle #1\rule[-1.4em]{0pt}{3.2em}$}
@@COURSEMACROS@@\newcommand{\UnitNumberName}{Unit @@UNITINT@@: @@UNITTITLE@@ \quad}
\newcommand{\LessonNumberName}{Lesson @@LESSONID@@: @@TITLE@@}

\begin{document}
\begin{center}
    {\Large\bfseries \CourseName \\
    {\small\bfseries \UnitNumberName \LessonNumberName}}
\end{center}

\begin{tcolorbox}[colback=forestbg, colframe=forest, boxrule=0.9pt, arc=2mm,
    left=2mm, right=2mm, top=1.5mm, bottom=1.5mm]
    \textbf{Primary Objective:} TODO --- what students can do/interpret/justify by the end.
    \\[2pt]
    \textbf{Standards (2023 VA SOL):} TODO --- the codes the user supplied, with clause notes.
    \\[2pt]
    \textbf{Lesson model:} \emph{Slides first.} The deck is the lesson and its printed handout is
    the student's notes. A warm-up on the slides; then, for each idea, a \textbf{definition}, a
    \textbf{worked example} (I do), and a \textbf{Now you try} worked alone in the handout's notes
    column before the answer is revealed on a click; a short wrap-up; and the \textbf{homework}
    started in class. \textbf{Homework source:} TODO --- generated in the packet / DeltaMath
    online (\emph{set name}) / DeltaMath printed.
\end{tcolorbox}

\renewcommand{\arraystretch}{1.4}
\begin{skillbox}[Priority Ideas \& Skills]{goldbox}
% Fill in the \item text ONLY. Two tabularx cells, one itemize each — do not re-nest or
% re-balance the environments (that is the classic source of a stray \end{itemize}).
\begin{tabularx}{\linewidth}{@{}X|X@{}}
\textbf{Priority Ideas \& Skills} & \textbf{Key Understandings (the ``why'')} \\[2pt]
\hline
\vspace{-6pt}
\begin{itemize}[leftmargin=1.1em, nosep, after=\vspace{2pt}]
  \item TODO: a skill, stated as something the student does.
\end{itemize}
&
\vspace{-6pt}
\begin{itemize}[leftmargin=1.1em, nosep, after=\vspace{2pt}]
  \item TODO: the why --- include the lesson's \textbf{target misconception} stated explicitly.
\end{itemize}
\end{tabularx}
\end{skillbox}

\begin{skillbox}[{Vocabulary, Concepts, \& Theorems --- one definition frame each}]{sky}
    % TODO: key terms + definitions, worded exactly as the deck's definition frames print them.
\end{skillbox}

\begin{fixedskillbox}[Lesson at a Glance --- \MeetingLength]{forestbg}
\renewcommand{\arraystretch}{1.25}
\begin{tabularx}{\linewidth}{@{}l c X X@{}}
\textbf{Phase} & \textbf{Min} & \textbf{Students} & \textbf{Teacher} \\
\hline
Warm-Up & 5 & TODO: the prior skills, alone, in the handout's notes column & Reveal the answers; land the question the first definition answers \\
Lesson & 40 & For each idea: read the definition, watch the example, then \emph{Now you try} alone before the reveal & TODO N cycles of definition $\to$ example $\to$ now you try; circulate during each try, then reveal \\
Wrap-up & 5 & Read back the definitions; the caution & The wrap-up frame; say the target misconception aloud \\
Homework & 10 & Start the homework alone & Launch item~1, then circulate \\
\end{tabularx}
\end{fixedskillbox}

\begin{skillbox}[Warm-Up (5 min, on the slides)]{forestbg}
    % TODO: the items, which prior skill each rehearses, and how the last one hands off into the
    % first definition.
\end{skillbox}

\begin{skillbox}[The Lesson --- definition, example, now you try (40 min)]{forestbg}
% One paragraph per cycle, in deck order: the definition (and what to point at on its display),
% the worked example and where students go wrong, the Now-you-try problem WITH its answer, and
% roughly how long to let them work before the reveal. Name the crux cycle.
\begin{multicols}{2}
\textbf{Cycle 1 --- TODO term.} TODO.

\textbf{Cycle 2 --- TODO.} TODO.

\columnbreak
\textbf{Cycle 3 --- TODO.} TODO --- normally the crux: the Now-you-try where the two answers
disagree; say so.
\end{multicols}
\end{skillbox}

\boxguard
\begin{skillbox}[Wrap-up (5 min)]{forestbg}
TODO: the definitions read back from the wrap-up frame, and the caution to say aloud.
\end{skillbox}

\boxguard
\begin{skillbox}[Homework --- scored, started in class, due the first class after two study halls]{goldbox}
\textbf{Source:} TODO --- \emph{generated} (the two packet pages) / \emph{DeltaMath online}
(the set name and what it covers) / \emph{DeltaMath printed} (the PDF dropped into the packet).
TODO: the items and what each covers; the formative check and the categories to sort responses
into. \textbf{Preview:} TODO next lesson.
\end{skillbox}

\boxguard
\begin{skillbox}[Active Monitoring --- Watch For]{redbox}
\begin{itemize}[leftmargin=1.2em, nosep]
  \item TODO: misconception, keyed to a Now-you-try or homework item number.
\end{itemize}
Cold-call prompts: TODO.
\end{skillbox}

% ── Teacher notes — one per component, in the order they are used ──
% This is the ONLY place teacher prose goes — never append one to a _key.
\begin{teachernote}[Slides]
TODO: the per-cycle pacing split that fills 40 minutes; which Now-you-try is the one that matters;
what to cut if you fall behind. Students write in the handout's notes column, which is printed
without the answers.
\end{teachernote}

\begin{teachernote}[Homework]
TODO: which item separates students and why; the formative check and its sort. The last 10
minutes are for \emph{starting} it in class --- launch item~1 aloud, then circulate; it is due the
first class after two study halls.
\end{teachernote}

\end{document}
